\documentclass[a4paper,12pt]{article} 
\usepackage[utf8]{inputenc} % Acentos válidos sin problemas
\usepackage[spanish]{babel} % Idioma
%-----------------------------------INICIO DE PACKETES------------------/
%----------------------------------------------------------------------/|
\usepackage{amsmath}   % Matemáticas: Comandos extras(cajas ecuaciones) |
\usepackage{amssymb}   % Matemáticas: Símbolos matemáticos              |
\usepackage{datetime}  % Agregar fechas                                 |
\usepackage{graphicx}  % Insertar Imágenes                              |
%\usepackage{biblatex} % Bibliografía                                   |
\usepackage{multicol}  % Creación de tablas                             |
\usepackage{longtable} % Tablas más largas                              |
\usepackage{xcolor}    % Permite cambiar colores del texto              |
\usepackage{tcolorbox} % Cajas de color                                 |
\usepackage{setspace}  % Usar espacios                                  |
\usepackage{fancyhdr}  % Para agregar encabezado y pie de página        |
\usepackage{float}     % Flotantes                                      |
\usepackage{soul}      % "Efectos" en palabras                          |
\usepackage{hyperref}  % Para usar hipervínculos                        |
\usepackage{caption}   % Utilizar las referencias                       |
\usepackage{subcaption} % Poder usar subfiguras                         |
\usepackage{multirow}  % Nos permite modificar tablas                   |
\usepackage{array}     % Permite utilizar los valores para multicolumn  |
\usepackage{booktabs}  % Permite modificar tablas                       |
\usepackage{diagbox}   % Diagonales para las tablas                     |
\usepackage{colortbl}  % Color para tablas                              |
\usepackage{listings}  % Escribir código                                |
\usepackage{mathtools} % SIMBOLO :=                                     |
\usepackage{enumitem}  % Modificar items de Listas                      |
\usepackage{tikz}      %                                                |
\usepackage{lipsum}    % for auto generating text                       |
\usepackage{afterpage} % for "\afterpage"s                              |
\usepackage{pagecolor} % With option pagecolor={somecolor or none}|     | 
%\usepackage{glossaries} %                                              |        
%----------------------------------------------------------------------\|
%-----------------------------------FIN--- DE PACKETES------------------\
%\usepackage{newtxtext,newtxmath}
%\usepackage{bm}
%--------------------------------/
%-------------------------------/
\usepackage[                 %   |
  headheight=15pt,  %            |
  letterpaper,  % Tipo de pag.   |
  left =1.5cm,  %  < 1 >         |
  right =1.5cm, %  < 1 >         | MARGENES DE LA PAGINA
  top =2cm,     %  < 1.5 >       |
  bottom =1.5cm %  < 1.5 >       |
]{geometry}     %                |
%-------------------------------\
%--------------------------------\

%----------------------------------------------------------------------/
%-------------------Encabezado y Pie de Pagina-----------------------/ |
%--------------------------------------------------------------------\ |
\usepackage{fancyhdr} %                                                |
\usepackage{lastpage} %                                                |
%\fancyhf{}           %                                                |
\pagestyle{fancy}     %                                                |
\fancyhead[L]{{$2024-1$}} %                                              |
\fancyhead[C]{}                      
\fancyhead[R]{}

\fancyfoot[L]{} 
\fancyfoot[C]{\thepage}
\fancyfoot[R]{} 

\renewcommand{\headrulewidth}{1.5pt} %                                 |
\renewcommand{\footrulewidth}{1.5pt} %                                 | 
%--------------------------------------------------------------------\ |
%-------------------Encabezado y Pie de Pagina-----------------------/ |
%----------------------------------------------------------------------/


\begin{document}%----------------------INICIO DOCUMENTO------------|
%------------------------------------------------------------------|

\newpage

\begin{titlepage}
    \center% Center everything on the page
    % Defines a new command for the horizontal lines, change thickness here
    \newcommand{\HRule}{\rule{\linewidth}{0.5mm}} 
    
    %---------------------
    %	LOGO
    %---------------------
    \includegraphics[width=4.5cm]{IMA/Ciencias.png} \\ 
    
    %----------------------------
    %	Encabezado
    %----------------------------
    \quad \\[1.5cm]
    %\textsc{\LARGE MSc Thesis}\\[1.5cm] % Name of your university/college
    \textsc{\Large Facultad de Ciencias}\\[0.5cm] % Major heading such as course name
    \textsc{\Large Computación Distribuida}\\[0.5cm] % Major heading such as course name
    %\textsc{\large Del Aula a la START UP}\\[0.5cm] % Minor heading such as course title
    %-------------
    %	TITULOS
    %-------------
    \makeatletter
        \HRule\\ [0.4cm]
            { \huge \bfseries PRACTICA 02}\\
        \HRule\\ [0.4cm]
    %-------------------
    %	PROFESORES
    %-------------------
    %\vspace{10mm}
    \large\textbf{Semestre $2024-1$}\\
    
    \vspace{10mm}
    \begin{minipage}{0.7\textwidth}
        \begin{flushleft} \large
            \emph{Profesor:}\\
                Luis Germán Pérez Hernández\\
                \vspace{3mm}
                \emph{Ayudantes:}\\
                Daniel Michel Tavera\\
                Yael Antonio Calzada Martín\\
        \end{flushleft}
    \end{minipage}
    %~
    \\
    %-------------------
    %	AUTOR
    %-------------------
    \vspace{10mm}
    \begin{minipage}{0.4\textwidth}
        \begin{flushright} \large
        \centering
            \emph{Alumnos} \\ % \emph{Alumnos:} \\
            \vspace{3mm}
            Marco Silva Huerta\\
            316205326\\
            \vspace{4mm}
            Edgar Montiel Ledesma\\
            317317794\\
            \vspace{4mm}
            Carlos Cortés\\
            420004846
            %\vspace{4mm}
            % Autor 2\\
        \end{flushright}
    \end{minipage}\\[1.2cm]
    \makeatother
    %---------------
    %	F E C H A
    %---------------
    {\large 18 de Octubre de 2023}\\[2cm]
    \vfill 
    \end{titlepage}

\newpage

\begin{center}
    {\huge Consenso}
\end{center}

\section*{Forma de compilar}

\section*{Funcionamiento}

\section*{Pseudocódigo del Algoritmo}

\subsection*{Algoritmo del Rey}
\begin{verbatim}
1. Definir las constantes:
   - NÚMERO_DE_GENERALES: número total de generales
   - NÚMERO_DE_TRAIDORES: número de generales traidores
   - F: número de generales traidores tolerados

2. Crear una estructura General con los siguientes campos:
   - id (entero): identificador del general
   - es_traidor (booleano): verdadero si el general es traidor, falso si es leal
   - voto (entero): 0 para retirada, 1 para ataque
   - mensaje (entero): mensaje enviado por el general en la ronda actual

3. Inicializar una lista de generales con NÚMERO_DE_GENERALES elementos.

4. Inicializar una variable REY con un valor aleatorio en el rango
               [0, NÚMERO_DE_GENERALES - 1]
   - Esto selecciona aleatoriamente a un general como el Rey sin 
     que los demás lo sepan.

5. Para cada general en la lista de generales:
   - Asignar un id único al general.
   - Determinar si el general es traidor (F generales serán traidores,
     incluyendo el Rey).
   - Inicializar el voto y el mensaje del general.

6. En cada ronda:
   - Cada general, incluido el Rey, elige su voto (0 para retirada, 1 para ataque) 
     de acuerdo a su estrategia.

7. Calcular el resultado de la ronda:
   - Inicializar las variables votos_ataque y votos_retirada a 0.
   - Para cada general en la lista de generales:
     - Si el general no es traidor:
       - Incrementar votos_ataque o votos_retirada según el voto del general.
     - Si el general es traidor:
       - Tomar el voto del general según su estrategia.

8. Verificar si la votación es válida:
   - Calcular la mayoría requerida como "(NÚMERO_DE_GENERALES / 2) + F".
   - Si votos_ataque >= mayoría o votos_retirada >= mayoría, la votación es válida.

9. Imprimir el resultado de la ronda y si la votación es válida o no.

10. Repetir las rondas hasta que se alcance un resultado válido o se llegue a un 
    límite de rondas.

11. Si se supera el límite de rondas, se considera que no hay consenso y se imprime
    un mensaje indicando la falta de consenso.

12. Finalizar el algoritmo.

\end{verbatim}

\section*{Desarrollo}








\end{document}%---------------------- FIN DOCUMENTO---------------|
%-----------------------------------------------------------------|